%%% FIELD def-hybrid-construction
Define
\[
 a^{\mathrm{loose}}_{nm}=
 \sqrt{\frac n{\mu r}}\left[1\wedge
 \frac B{\lambda_{nm}}\left\{
 \sqrt{\frac{(n\vee m)\log(e(n+m))}{q}}+
 \frac{\log(e(n+m))}{q}\right\}\right].
\]
For \(i\in I\) and \(s\in\{1,\ldots,r\}\), set
\[
 U^{\mathrm D}_{is}=\sqrt{\frac{2}{n+1}}
 \sin\!\left(\frac{\pi i s}{n+1}\right).
\]
The hybrid row-factor estimator is
\[
 \widehat U^{\mathrm H}=
 \begin{cases}
 \widehat U^{\mathrm I},&a^{\mathrm{loose}}_{nm}\le1,\\
 U^{\mathrm D},&a^{\mathrm{loose}}_{nm}>1.
 \end{cases}
\]

%%% FIELD lem-sine-statement
The deterministic discrete-sine matrix \(U^{\mathrm D}\) belongs to \(\mathbb O_{n,r}\), satisfies \(\lVert U^{\mathrm D}\rVert_{2,\infty}\le\sqrt{2r/n}\), and, for every \(U\in\mathbb O_{n,r}\) satisfying \(\lVert U\rVert_{2,\infty}\le\sqrt{\mu r/n}\), obeys the pathwise bound
\[
 \mathcal L(U^{\mathrm D},U)\le1+\sqrt{2/\mu}.
\]

%%% FIELD lem-sine-proof
For \(h\in\{1,\ldots,2n\}\), put
\[
 C_h=\sum_{i=1}^n\cos\!\left(\frac{\pi h i}{n+1}\right),
 \qquad C_0=n.
\]
If \(h\) is odd, the finite-cosine-sum identity, obtained by taking real parts of the finite geometric series, gives
\[
 C_h=
 \frac{\sin\{n\pi h/(2(n+1))\}\cos(\pi h/2)}
      {\sin\{\pi h/(2(n+1))\}}=0.
\]
Every denominator is nonzero because \(1\le h\le2n<2(n+1)\).  If \(h=2a\) is even, where \(1\le a\le n\), then
\[
 1+\sum_{i=1}^n\exp\!\left(\frac{2\pi\mathrm i a i}{n+1}\right)=0
\]
by the geometric-series formula, because its ratio is a nontrivial \((n+1)\)-st root of unity.  Taking real parts yields \(C_h=-1\).

For \(k,s\in\{1,\ldots,r\}\), the product-to-sum identity therefore gives
\[
 \begin{aligned}
 \sum_{i=1}^n
 \sin\!\left(\frac{\pi i k}{n+1}\right)
 \sin\!\left(\frac{\pi i s}{n+1}\right)
 &=\frac12\{C_{|k-s|}-C_{k+s}\}\\
 &=\frac{n+1}{2}\,\mathbf1\{k=s\}.
 \end{aligned}
\]
Indeed, when \(k\ne s\), the two positive indices \(|k-s|\) and \(k+s\) have the same parity and their corresponding \(C\)'s agree; when \(k=s\), \(C_0=n\) and \(C_{2k}=-1\).  The definition in \(\mathrm{def:hybrid\mbox{-}row\mbox{-}estimator}\) now implies
\[
 (U^{\mathrm D})^\top U^{\mathrm D}=I_r,
\]
where \(r\le n\) by the declared range of \(r\).  Thus \(U^{\mathrm D}\in\mathbb O_{n,r}\).

For every row \(i\), \(|\sin x|\le1\) gives
\[
 \lVert e_i^\top U^{\mathrm D}\rVert_2^2
 =\frac2{n+1}\sum_{s=1}^r
 \sin^2\!\left(\frac{\pi i s}{n+1}\right)
 \le\frac{2r}{n+1}\le\frac{2r}{n}.
\]
Taking the maximum over rows proves the asserted incoherence bound.  Finally, choose \(R=I_r\) in the minimum defining \(\mathcal L\).  The triangle inequality for every row, the preceding display, and the assumed left-incoherence bound give
\[
 \begin{aligned}
 \mathcal L(U^{\mathrm D},U)
 &\le\sqrt{\frac n{\mu r}}
 \lVert U^{\mathrm D}-U\rVert_{2,\infty}\\
 &\le\sqrt{\frac n{\mu r}}
 \left(\sqrt{\frac{2r}{n}}+\sqrt{\frac{\mu r}{n}}\right)
 =1+\sqrt{\frac2\mu}.
 \end{aligned}
\]
This is deterministic and holds for every admissible target \(U\).

%%% FIELD prop-hybrid-statement
There is a universal constant \(C<\infty\) such that, for every \(n,m\ge2\), the hybrid estimator satisfies \(\widehat U^{\mathrm H}\in\mathfrak E_{n,m}\) and
\[
 \sup_{P\in\mathfrak P^{\mathrm{pred}}_{n,m}(r,\mu,K,q,B,\lambda)}
 E_P\!\left[\mathcal L(\widehat U^{\mathrm H},U)\right]
 \le C\bigl(1\wedge a^{\mathrm{loose}}_{nm}\bigr).
\]
When \(a^{\mathrm{loose}}_{nm}>1\), the bound is delivered by the deterministic discrete-sine frame and is pathwise; when \(a^{\mathrm{loose}}_{nm}\le1\), it is the recorded IPW--SVD expected-risk bound.

%%% FIELD prop-hybrid-proof
The score \(Z^{\mathrm{ipw}}\) is a measurable function of \((A,Y,p)\) by \(\mathrm{def:ipw\mbox{-}score}\).  Rank projection followed by the fixed tie rule is a Borel-measurable finite-dimensional map, so \(\widehat U^{\mathrm I}\) is \(\mathcal O_{n,m}\)-measurable by \(\mathrm{def:ipw\mbox{-}svd}\) and \(\mathrm{def:observed\mbox{-}estimators}\).  Lemma \(\mathrm{lem:discrete\mbox{-}sine\mbox{-}fallback\mbox{-}loss}\) shows that \(U^{\mathrm D}\in\mathbb O_{n,r}\); being deterministic, it is also \(\mathcal O_{n,m}\)-measurable.  The event \(\{a^{\mathrm{loose}}_{nm}\le1\}\) is deterministic because the radius in \(\mathrm{def:hybrid\mbox{-}row\mbox{-}estimator}\) depends only on the declared class parameters and dimensions.  Both branches of \(\widehat U^{\mathrm H}\) are therefore \(\mathcal O_{n,m}\)-measurable and take values in \(\mathbb O_{n,r}\).  Hence
\[
 \widehat U^{\mathrm H}\in\mathfrak E_{n,m}.
\]

Fix \(P\in\mathfrak P^{\mathrm{pred}}_{n,m}(r,\mu,K,q,B,\lambda)\).  If \(a^{\mathrm{loose}}_{nm}\le1\), then \(\widehat U^{\mathrm H}=\widehat U^{\mathrm I}\), and Proposition \(\mathrm{prop:ipw\mbox{-}expected\mbox{-}risk\mbox{-}upper}\) gives
\[
 E_P\mathcal L(\widehat U^{\mathrm H},U)
 \le C_0a^{\mathrm{loose}}_{nm}
 =C_0(1\wedge a^{\mathrm{loose}}_{nm}).
\]
If \(a^{\mathrm{loose}}_{nm}>1\), then \(\widehat U^{\mathrm H}=U^{\mathrm D}\).  Membership of \(P\) in the predictable class includes \(\mathrm{ass:left\mbox{-}incoherence}\), so Lemma \(\mathrm{lem:discrete\mbox{-}sine\mbox{-}fallback\mbox{-}loss}\) yields, pathwise,
\[
 \mathcal L(\widehat U^{\mathrm H},U)
 \le1+\sqrt{2/\mu}\le1+\sqrt2,
\]
where the last inequality uses \(\mu\ge1\).  Since \(1\wedge a^{\mathrm{loose}}_{nm}=1\) in this case, taking \(C=\max(C_0,1+\sqrt2)\), then taking expectations and the supremum over \(P\), proves the proposition.

%%% FIELD prop-hybrid-bracket-statement
Fix \(r\), \(K\), and \(\mu>1\).  There are constants \(c_0,c_1,c,C>0\), depending only on these fixed quantities, \(q\), and the aspect bounds, such that, for all sufficiently large \(n,m\), whenever \(b_{nm}\le c_0\) and \(\lambda_{nm}\le c_1B\sqrt{nm}\),
\[
 c b_{nm}\le
 \inf_{\widetilde U\in\mathfrak E_{n,m}}
 \sup_{P\in\mathfrak P^{\mathrm{pred}}_{n,m}(r,\mu,K,q,B,\lambda)}
 E_P[\mathcal L(\widetilde U,U)]
 \le
 \sup_{P\in\mathfrak P^{\mathrm{pred}}_{n,m}(r,\mu,K,q,B,\lambda)}
 E_P[\mathcal L(\widehat U^{\mathrm H},U)]
 \le C\bigl(1\wedge a^{\mathrm{loose}}_{nm}\bigr).
\]
The lower endpoint is the constant-propensity information bound and the upper endpoint is attained by one observed-data hybrid estimator.  The separate guarantee \(E_P\mathcal L(\widehat U^{\mathrm I},U)\le C a^{\mathrm{loose}}_{nm}\) remains valid, while the endogenous-target obstruction concerns only recorded IPW--SVD and proves no constant lower bound for all estimators.

%%% FIELD prop-hybrid-bracket-proof
The first inequality is Theorem \(\mathrm{thm:adaptive\mbox{-}minimax\mbox{-}lower}\), applied under its displayed conditions to the same predictable class and expected normalized maximum-row loss.  Proposition \(\mathrm{prop:hybrid\mbox{-}expected\mbox{-}risk\mbox{-}upper}\) proves both that \(\widehat U^{\mathrm H}\in\mathfrak E_{n,m}\) and that its worst-case expected loss is at most \(C(1\wedge a^{\mathrm{loose}}_{nm})\).  Evaluating the infimum over \(\mathfrak E_{n,m}\) at this one admissible estimator gives the middle inequality, and the same proposition gives the last inequality.

Proposition \(\mathrm{prop:ipw\mbox{-}expected\mbox{-}risk\mbox{-}upper}\) separately gives the stated \(Ca^{\mathrm{loose}}_{nm}\) risk bound for \(\widehat U^{\mathrm I}\).  Proposition \(\mathrm{prop:adaptive\mbox{-}endogenous\mbox{-}target\mbox{-}obstruction}\) evaluates that particular estimator on one sequence; it does not lower-bound the risk of \(\widehat U^{\mathrm H}\) or the infimum over all of \(\mathfrak E_{n,m}\).  This proves the bracket and the stated logical separation.

%%% FIELD thm-lower-proof
Put \(t=m/n\).  Assumption \(\mathrm{ass:aspect\mbox{-}ratio}\) gives \(t\in[c_-,c_+]\).  Choose \(c_0<1\).  On the regime \(b_{nm}\le c_0\), the truncation in \(\mathrm{def:frontier\mbox{-}rate}\) is inactive, and hence
\[
 \begin{aligned}
 b_{nm}
 &=\frac{B}{\lambda_{nm}}K r^{3/2}\mu q^{-1/2}
   \sqrt{\log(en)}\sqrt{nm}\sqrt{n^{-1}+nm^{-2}}\\
 &=K r^{3/2}\mu\sqrt{t+t^{-1}}
   \frac{B}{\lambda_{nm}}\sqrt{\frac{n\log(en)}q}.
 \end{aligned}
\]
Thus, with
\[
 s_{nm}:=\frac{B}{\lambda_{nm}}\sqrt{\frac{n\log(en)}q},
\]
there are positive constants \(C_-\) and \(C_+\), depending only on \(r,K,\mu,c_-,c_+\), such that
\[
 C_-s_{nm}\le b_{nm}\le C_+s_{nm}.
 \tag{1}
\]

We construct a constant-propensity subexperiment.  Since \(r\) is fixed, for all sufficiently large dimensions partition \(I\) into nonempty disjoint balanced blocks \(I_1,\ldots,I_r\), and partition \(J\) into nonempty disjoint balanced blocks \(J_1,\ldots,J_r\).  Write \(N_s=|I_s|\) and \(M_s=|J_s|\).  For \(s\ge2\), let \(u_s=N_s^{-1/2}\mathbf1_{I_s}\), viewed in \(\mathbb R^n\), and, for every \(s\), let \(v_s=M_s^{-1/2}\mathbf1_{J_s}\), viewed in \(\mathbb R^m\).  Let
\[
 u_0=N_1^{-1/2}\mathbf1_{I_1}.
\]
Fix \(k_0\in I_1\).  For every \(k\in\mathcal K:=I_1\setminus\{k_0\}\), define
\[
 w_k=\frac{e_k-e_{k_0}}{\sqrt2},\qquad
 \delta=a\frac{B}{\lambda_{nm}}\sqrt{\frac{\log(en)}q},\qquad
 u_k=\frac{u_0+\delta w_k}{\sqrt{1+\delta^2}},
 \tag{2}
\]
where the fixed constant \(a>0\) will be chosen below.  The vectors \(u_0\) and \(w_k\) are orthogonal unit vectors.  Hence \(u_k\) is a unit vector.  Disjoint block supports show that
\[
 U_k=(u_k,u_2,\ldots,u_r)\in\mathbb O_{n,r},
 \qquad V=(v_1,\ldots,v_r)\in\mathbb O_{m,r}.
\]
Also set \(U_0=(u_0,u_2,\ldots,u_r)\), and define
\[
 M_k=\lambda_{nm}U_kV^\top\quad(k\in\mathcal K),
 \qquad M_0=\lambda_{nm}U_0V^\top.
 \tag{3}
\]

We verify the structural restrictions.  From (1)--(2),
\[
 \delta\sqrt n=a s_{nm}\le aC_-^{-1}b_{nm}\le aC_-^{-1}c_0.
 \tag{4}
\]
For a row in \(I_1\), the triangle inequality gives
\[
 \lVert e_i^\top U_k\rVert_2
 \le N_1^{-1/2}+\delta/\sqrt2.
 \tag{5}
\]
For a row in \(I_s\), \(s\ge2\), its row norm is \(N_s^{-1/2}\).  Balanced blocks give
\[
 \max_s\sqrt{n/N_s}\longrightarrow\sqrt r.
\]
The strict inequality \(\mu>1\) supplies positive slack between \(\sqrt r\) and \(\sqrt{\mu r}\).  First choose \(a c_0\) small enough, depending only on that slack and the constants in (4), and then take \(n\) sufficiently large.  Equations (4)--(5) yield
\[
 \lVert U_k\rVert_{2,\infty}\le\sqrt{\mu r/n}.
 \tag{6}
\]
The same balanced-block calculation, without the perturbation, gives
\[
 \lVert V\rVert_{2,\infty}\le\sqrt{\mu r/m}
 \tag{7}
\]
for all sufficiently large \(m\).  By (3), all \(r\) nonzero singular values of \(M_k\) equal \(\lambda_{nm}\).  Thus exact rank, the signal floor, and the condition-number bound hold, the latter because \(1\le K\).

It remains to make the potential outcomes bounded.  For all sufficiently large dimensions, balanced blocks satisfy \(N_s\ge n/(2r)\) and \(M_s\ge m/(2r)\).  Equations (2)--(3) imply
\[
 \max_{i,j}|(M_k)_{ij}|
 \le \frac{2r\lambda_{nm}}{\sqrt{nm}}
     +aB\sqrt{\frac{2r\log(en)}{qm}}.
 \tag{8}
\]
Choose \(c_1>0\) so that \(2rc_1\le1/4\).  Under \(\lambda_{nm}\le c_1B\sqrt{nm}\), the first term in (8) is at most \(B/4\).  The second is at most \(B/4\) for all sufficiently large \(n,m\), because \(q,r,a\) are fixed and the aspect ratio is bounded.  Hence
\[
 \max_{i,j}|(M_k)_{ij}|\le B/2
 \tag{9}
\]
uniformly over \(k\in\mathcal K\cup\{0\}\).

For each \(k\in\mathcal K\cup\{0\}\), define a law \(P_k\) as follows.  Set \(Y_{ij}(0)=0\).  Independently over cells, draw \(Y_{ij}(1)\in\{-B,B\}\) with
\[
 P_k\{Y_{ij}(1)=B\}
 =h_{k,ij}:=\frac12\left(1+\frac{(M_k)_{ij}}B\right).
 \tag{10}
\]
Independently of all potential outcomes and over cells, draw \(A_{ij}\sim\operatorname{Bernoulli}(q)\), and record \(p_{ij}=q\).  Equation (9) puts every \(h_{k,ij}\) in \([1/4,3/4]\), while (10) gives \(E_{P_k}[Y_{ij}(1)]=(M_k)_{ij}\).  Independence from all preceding cells in the filtration of \(\mathrm{def:global\mbox{-}filtration}\) gives
\[
 \mu^1_{ij}=E_{P_k}[Y_{ij}(1)\mid\mathbb G_{o(i,j)-1}]=(M_k)_{ij},
 \qquad \mu^0_{ij}=0.
 \tag{11}
\]
Thus the causal effect matrix \(M=\mu^1-\mu^0\) is exactly \(M_k\).  Independence of \(A_{ij}\) from the potential outcomes, also conditional on the recorded past, verifies sequential randomization.  The equality \(p_{ij}=q\) verifies overlap because \(q<1/2\), and (10), together with \(Y_{ij}(0)=0\), verifies bounded outcomes.  Equations (6)--(7), the singular-value calculation, and the aspect-ratio assumption verify every remaining member property.  Therefore
\[
 P_k\in\mathfrak P^{\mathrm{pred}}_{n,m}(r,\mu,K,q,B,\lambda)
 \quad\text{for every }k\in\mathcal K\cup\{0\}.
 \tag{12}
\]

We next bound the information in \(\mathcal O_{n,m}=\sigma(A,Y,p)\).  If \(x,y\in[-B/2,B/2]\), put \(s=(1+x/B)/2\) and \(t'=(1+y/B)/2\).  For fixed \(t'\in[1/4,3/4]\), Bernoulli relative entropy \(f(s)=s\log(s/t')+(1-s)\log((1-s)/(1-t'))\) satisfies
\[
 f(t')=f'(t')=0,
 \qquad f''(s)=\frac1{s(1-s)}\le\frac{16}3.
\]
Taylor's theorem gives
\[
 \operatorname{kl}\left(\frac{1+x/B}{2},\frac{1+y/B}{2}\right)
 \le \frac{2}{3}\frac{(x-y)^2}{B^2}.
 \tag{13}
\]
The assignment law and recorded propensity are identical under \(P_k\) and \(P_0\).  Conditional on \(A_{ij}=0\), the observed outcome is the common value zero; conditional on \(A_{ij}=1\), it has Bernoulli sign parameter \(h_{k,ij}\).  Independence over cells, (13), and direct factorization of the product likelihood give
\[
 \begin{aligned}
 \operatorname{KL}(P_k^{\mathcal O},P_0^{\mathcal O})
 &=q\sum_{i,j}\operatorname{kl}(h_{k,ij},h_{0,ij})\\
 &\le \frac{Cq}{B^2}\lVert M_k-M_0\rVert_F^2.
 \end{aligned}
 \tag{14}
\]
Only the first singular component differs, so
\[
 \lVert M_k-M_0\rVert_F^2
 =\lambda_{nm}^2\lVert u_k-u_0\rVert_2^2.
\]
Orthogonality in (2) gives
\[
 \lVert u_k-u_0\rVert_2^2
 =2\left(1-\frac1{\sqrt{1+\delta^2}}\right)
 \le\delta^2.
\]
Substitution into (14) proves
\[
 \operatorname{KL}(P_k^{\mathcal O},P_0^{\mathcal O})
 \le Ca^2\log(en).
 \tag{15}
\]

The alternatives are separated in rowwise loss.  If \(k,\ell\in\mathcal K\) are distinct, inspect row \(k\).  On that row only the first columns of \(U_k\) and \(U_\ell\) can be nonzero, and (2) gives
\[
 \lVert e_k^\top U_k\rVert_2-\lVert e_k^\top U_\ell\rVert_2
 =\frac{\delta}{\sqrt{2(1+\delta^2)}}.
\]
Right multiplication by any \(R\in\mathbb O_r\) preserves a row norm.  The reverse triangle inequality therefore gives
\[
 \min_{R\in\mathbb O_r}\lVert U_kR-U_\ell\rVert_{2,\infty}
 \ge\frac{\delta}{\sqrt{2(1+\delta^2)}}.
\]
By \(\mathrm{def:rowwise\mbox{-}loss}\),
\[
 \mathcal L(U_k,U_\ell)
 \ge \Delta_{nm}:=\sqrt{\frac n{\mu r}}
       \frac{\delta}{\sqrt{2(1+\delta^2)}}.
 \tag{16}
\]
Equation (4) implies \(\delta\to0\), so eventually \(\delta\le1\).  Combining (1), (2), and (16), with \(t\in[c_-,c_+]\), gives
\[
 \Delta_{nm}\ge c b_{nm}.
 \tag{17}
\]
Indeed, before the bounded denominator factor in (16),
\[
 \frac{\sqrt{n/(\mu r)}\,\delta}{b_{nm}}
 =\frac{a}{K r^2\mu^{3/2}\sqrt{t+t^{-1}}}.
\]

We derive the multiple-testing reduction.  Let \(D=|\mathcal K|=N_1-1\), let \(\Theta\) be uniform on \(\mathcal K\), and, conditional on \(\Theta=k\), draw the observed experiment from \(P_k^{\mathcal O}\).  With \(\overline P=D^{-1}\sum_{k\in\mathcal K}P_k^{\mathcal O}\), expansion of logarithms gives
\[
 \begin{aligned}
 I(\Theta;\mathcal O)
 &=\frac1D\sum_{k\in\mathcal K}\operatorname{KL}(P_k^{\mathcal O},\overline P)\\
 &=\frac1D\sum_{k\in\mathcal K}\operatorname{KL}(P_k^{\mathcal O},P_0^{\mathcal O})
  -\operatorname{KL}(\overline P,P_0^{\mathcal O})\\
 &\le Ca^2\log(en),
 \end{aligned}
 \tag{18}
\]
where the last line uses nonnegativity of relative entropy and (15).  For any decoder \(\widehat\Theta\), let \(P_e=P(\widehat\Theta\ne\Theta)\).  Conditioning the entropy of \(\Theta\) first on the error indicator yields
\[
 H(\Theta\mid\mathcal O)\le\log2+P_e\log D.
\]
Since \(H(\Theta)=\log D\), (18) implies
\[
 P_e\ge1-\frac{Ca^2\log(en)+\log2}{\log D}.
 \tag{19}
\]
Balanced blocks give \(D\asymp n\).  Choose \(a>0\) sufficiently small, depending only on the fixed quantities, so that the right side of (19) is at least \(1/2\) for all sufficiently large \(n\).

Finally take any \(\widetilde U\in\mathfrak E_{n,m}\) and decode by a nearest member of \(\{U_k:k\in\mathcal K\}\) in the orthogonal-quotient \(2,\infty\) distance.  This quotient distance obeys the triangle inequality: compose two orthogonal alignments and use invariance of row norms under right orthogonal multiplication.  On a decoding error under \(P_k\), pairwise separation (16) and nearest-neighbor decoding imply
\[
 \mathcal L(\widetilde U,U_k)\ge\Delta_{nm}/2.
\]
Averaging over the uniform index, using (19), and then (17), gives
\[
 \sup_{k\in\mathcal K}E_{P_k}[\mathcal L(\widetilde U,U_k)]
 \ge\frac{\Delta_{nm}}4
 \ge c_2b_{nm}.
\]
By (12), the finite family is a subset of the predictable class and has \(p_{ij}=q\) at every cell.  Taking the supremum over the full class and then the infimum over all \(\widetilde U\in\mathfrak E_{n,m}\) proves the theorem.

%%% FIELD thm-ipw-proof
Proposition \(\mathrm{prop:ipw\mbox{-}expected\mbox{-}risk\mbox{-}upper}\) gives the stronger uniform first-moment statement
\[
 \sup_{P\in\mathfrak P^{\mathrm{pred}}_{n,m}(r,\mu,K,q,B,\lambda)}
 E_P\!\left[\mathcal L(\widehat U^{\mathrm I},U)\right]
 \le C a^{\mathrm{loose}}_{nm}.
\]
The radius is positive because \(B,q,\lambda_{nm}>0\) and \(n,m\ge2\).  Markov's inequality therefore gives, for every \(t>0\),
\[
 \sup_{P\in\mathfrak P^{\mathrm{pred}}_{n,m}(r,\mu,K,q,B,\lambda)}
 P\!\left\{\mathcal L(\widehat U^{\mathrm I},U)>
 t a^{\mathrm{loose}}_{nm}\right\}\le \frac Ct.
\]
This is the asserted uniform \(O_P(a^{\mathrm{loose}}_{nm})\) bound.  The radius and the moment proof use only the score matrix and contain no query loading, so no factor depending on \(|\mathcal G_n|\) occurs.

For the nonreplacement assertion, Proposition \(\mathrm{prop:adaptive\mbox{-}endogenous\mbox{-}target\mbox{-}obstruction}\) supplies a sequence \(P_n\) in the same class and constants \(c_0,c_1>0\) such that \(b_{nn}\to0\) and
\[
 \liminf_{n\to\infty}P_n\!\left\{
 \mathcal L(\widehat U^{\mathrm I},U)>c_0\right\}\ge c_1.
\]
If the IPW--SVD loss were uniformly \(O_P(b_{nm})\), then, for \(\varepsilon=c_1/2\), some fixed \(T<\infty\) would make the probability at threshold \(Tb_{nn}\) at most \(\varepsilon\) eventually.  Since \(Tb_{nn}<c_0\) eventually, this contradicts the preceding display.  The witness evaluates only \(\widehat U^{\mathrm I}\), so it proves no impossibility result for a different observed-data estimator.

%%% FIELD prop-ipw-proof
Fix \(P\) in the predictable class and expose cells in the row-major order of \(\mathrm{def:global\mbox{-}filtration}\).  By Lemma \(\mathrm{lem:ipw\mbox{-}identification}\), each \(E_{ij}\) has conditional mean zero given \(\mathbb G_{o(i,j)-1}\).  Sequential randomization, bounded outcomes, and overlap give
\[
 \begin{aligned}
 E_P\!\left[(Z^{\mathrm{ipw}}_{ij})^2
 \mid\mathbb G_{o(i,j)-1}\right]
 &=E_P\!\left[
 \frac{Y_{ij}(1)^2}{p_{ij}}+
 \frac{Y_{ij}(0)^2}{1-p_{ij}}
 \mathrel{\Big|}\mathbb G_{o(i,j)-1}\right]\\
 &\le \frac{2B^2}{q}.
 \end{aligned}
 \tag{20}
\]
For the equality, condition first on the two potential outcomes and use \(\mathrm{ass:sequential\mbox{-}randomization}\); both divisions are valid by \(\mathrm{ass:overlap}\).  Since \(M_{ij}=E_P[Z^{\mathrm{ipw}}_{ij}\mid\mathbb G_{o(i,j)-1}]\), conditional variance is no larger than (20).  Moreover, \(|Z^{\mathrm{ipw}}_{ij}|\le B/q\), while conditional Jensen and bounded outcomes give \(|M_{ij}|\le2B\).  Since \(q<1/2\),
\[
 |E_{ij}|\le \frac{2B}{q}.
 \tag{21}
\]

Introduce the self-adjoint dilation
\[
 \Delta_{ij}=
 \begin{pmatrix}
 0&E_{ij}e_ie_j^\top\\
 E_{ij}e_je_i^\top&0
 \end{pmatrix}.
\]
Its norm is \(|E_{ij}|\), and its square has diagonal blocks \(E_{ij}^2e_ie_i^\top\) and \(E_{ij}^2e_je_j^\top\).  Hence (20) gives
\[
 \left\lVert\sum_{i,j}E_P[\Delta_{ij}^2
 \mid\mathbb G_{o(i,j)-1}]\right\rVert
 \le v_{nm}:=\frac{2B^2(n\vee m)}q.
\]
Lemma \(\mathrm{lem:tropp\mbox{-}matrix\mbox{-}freedman}\), applied to both signs of this \((n+m)\)-dimensional martingale and using (21), gives
\[
 P\{\lVert E\rVert\ge t\}
 \le2(n+m)\exp\!\left\{-\frac{t^2/2}{v_{nm}+R_{nm}t/3}\right\},
 \qquad R_{nm}=\frac{2B}{q}.
 \tag{22}
\]
The terminal dilation has operator norm \(\lVert E\rVert\).  Put \(L_{nm}=\log(2(n+m))\).  Solving the quadratic inequality in (22) shows that, for every \(x\ge0\),
\[
 P\!\left\{\lVert E\rVert>
 \sqrt{2v_{nm}(L_{nm}+x)}+
 \frac{2R_{nm}}3(L_{nm}+x)\right\}\le e^{-x}.
 \tag{23}
\]
To integrate (23), let \(T(x)\) denote its displayed threshold.  It is increasing, so the layer-cake formula and the change of variables \(t=T(x)\) give
\[
 E_P\lVert E\rVert
 \le T(0)+\int_0^\infty e^{-x}T'(x)\,dx.
\]
Because \(L_{nm}\ge1\), the integral is bounded by a universal constant times \(\sqrt{v_{nm}L_{nm}}+R_{nm}L_{nm}\).  Consequently
\[
 E_P\lVert E\rVert
 \le C B\left\{
 \sqrt{\frac{(n\vee m)\log(e(n+m))}{q}}+
 \frac{\log(e(n+m))}{q}\right\}
 =:C\eta_{nm}.
 \tag{24}
\]
Every bound through (24) is uniform in \(P\).

It remains to pass from operator error to the declared loss.  Let \(P_U=UU^\top\), and let \((\widehat U,\widehat\Sigma,\widehat V)\) be the leading rank-\(r\) singular triple of \(M+E\).  Exact rank gives \((I-P_U)M=0\).  The signal floor and the singular-value min--max formula imply, on \(\{\lVert E\rVert\le\lambda_{nm}/2\}\),
\[
 \sigma_r(\widehat\Sigma)
 \ge\sigma_r(M)-\lVert E\rVert
 \ge\lambda_{nm}-\lVert E\rVert>0.
\]
Projecting \((M+E)\widehat V=\widehat U\widehat\Sigma\) onto the orthogonal complement of \(U\) gives
\[
 \lVert(I-P_U)\widehat U\rVert
 \le\frac{\lVert E\rVert}{\lambda_{nm}-\lVert E\rVert}
 \le\frac{2\lVert E\rVert}{\lambda_{nm}}.
\]
Diagonalizing the principal-angle cosines of \(U^\top\widehat U\) and using \(2(1-\cos t)\le2\sin^2t\) supplies an \(R\in\mathbb O_r\) such that
\[
 \lVert\widehat UR-U\rVert
 \le\sqrt2\lVert(I-P_U)\widehat U\rVert.
\]
Off the event \(\{\lVert E\rVert\le\lambda_{nm}/2\}\), the aligned operator distance is at most two.  Since \(\lVert Q\rVert_{2,\infty}\le\lVert Q\rVert\), the two cases combine into
\[
 \mathcal L(\widehat U^{\mathrm I},U)
 \le C\sqrt{\frac n{\mu r}}
 \left(1\wedge\frac{\lVert E\rVert}{\lambda_{nm}}\right).
 \tag{25}
\]
The fixed tie rule ensures that the leading left singular vectors of \(\mathcal P_r(M+E)\) are the chosen leading left singular vectors of \(M+E\).  Finally, \(x\mapsto1\wedge x\) is concave on \(\mathbb R_+\).  Jensen's inequality, (24), and (25) yield
\[
 \begin{aligned}
 E_P\mathcal L(\widehat U^{\mathrm I},U)
 &\le C\sqrt{\frac n{\mu r}}
 E_P\left[1\wedge\frac{\lVert E\rVert}{\lambda_{nm}}\right]\\
 &\le C\sqrt{\frac n{\mu r}}
 \left(1\wedge\frac{\eta_{nm}}{\lambda_{nm}}\right)
 \le C a^{\mathrm{loose}}_{nm}.
 \end{aligned}
\]
Taking the supremum over the declared class proves the proposition.

%%% FIELD oeq-proposed-statement
Along the explicit path \(n_\ell=2\ell\) and \(m_\ell=2\lceil\gamma\ell\rceil\), write
\[
 R^\star_\ell=
 \inf_{\widetilde U\in\mathfrak E_{n_\ell,m_\ell}}
 \sup_{P\in\mathfrak P^{\mathrm{pred}}_{n_\ell,m_\ell}(r,\mu,K,q,B,\lambda)}
 E_P[\mathcal L(\widetilde U,U)].
\]
In the nonempty lower-bound regime with \(\mu>1\) and \(b_{n_\ell m_\ell}\to0\), determine whether \(R^\star_\ell\) is of order \(b_{n_\ell m_\ell}\), has a strictly intermediate order between \(b_{n_\ell m_\ell}\) and a constant, or remains bounded away from zero.  Only after resolving this rate question, determine whether \(C_\star\) is finite and, when finite, its exact value and an attaining observed-data estimator.

%%% FIELD oeq-partial-result
The proved common-risk bracket is
\[
 c b_{nm}\le
 \inf_{\widetilde U\in\mathfrak E_{n,m}}
 \sup_{P\in\mathfrak P^{\mathrm{pred}}_{n,m}(r,\mu,K,q,B,\lambda)}
 E_P\mathcal L(\widetilde U,U)
 \le C(1\wedge a^{\mathrm{loose}}_{nm}).
\]
The constant-propensity Fano family proves the lower endpoint.  The deterministic-frame hybrid proves the upper endpoint with one observed-data estimator.  On every nonempty bounded-outcome class, \(\lambda_{nm}\le2B\sqrt{nm}\); under fixed \(r,\mu,K,q\) and bounded aspect ratio, this forces \(b_{nm}\gtrsim\sqrt{\log n/n}\) and \(a^{\mathrm{loose}}_{nm}\gtrsim\sqrt{\log n}\).  Thus the available uniform upper resolution selects the constant fallback in that regime.  The endogenous-target witness rules out \(b_{nm}\)-attainment only for recorded IPW--SVD, not for every estimator.
